% Teaser: title block plus the full-width figure above the abstract.
% Replaces \maketitle in main.tex. Figure sources are in assets/teaser/.
\twocolumn[{%
\renewcommand\twocolumn[1][]{#1}%
\maketitle
\begin{center}
    \captionsetup{type=figure}
    \vspace{-2mm}
    \includegraphics[width=\textwidth]{assets/teaser.png}
    \captionof{figure}{\textbf{A bimanual robot is two arms, not one arm with
    twice the joints.}
    Left: existing tokenizers treat the two arms as a single arm and mix the
    \textcolor[HTML]{1F6676}{\textbf{left}} and
    \textcolor[HTML]{C96A2B}{\textbf{right}} arm in every token; BiCAT gives
    each arm its own token stream, aware of its partner.
    Right: BiCAT rollouts on three embodiments with each arm's motion over
    time, and success rates in simulation and on real robots; white numbers
    give the gain over the best baseline.}
    % NOTE (figure edit, not caption text): rename "Tokenizer" to "Action Tokenizer" in teaser.pdf.
    % Reviewer note from Overleaf (2026-11-08), addressed by the caption above:
    % The figure is good, but i think it doesn't tell me what is your contribution,
    % like you want to show that most tokenizer is optimize of design for single arm,
    % but not for bimanual, and why yours is design for bimanual?
    \label{fig:teaser}
\end{center}%
\vspace{2mm}
}]
